\chapter{Learning Systems and Algorithmic Histories}
\index{learning systems}\index{polynomial sketch}

\section{The Reconstruction Problem in Learning Systems}

A machine-learning model is trained on a dataset of examples. When training ends, the parameters are fixed and the training history is discarded. The counterfactual question: what would the model have predicted had a particular training example been excluded?

Retraining large models for every such question is computationally prohibitive. The question is whether the model's current parameters contain enough information to answer such counterfactual questions without retraining.

\section{Interaction Decomposition: Diagnostic Question 1}

Let $\mathcal{X} = \mathbb{R}^n$, where $x_i$ is a downweight on the $i$-th training example. The query is $Q(x) = g(f(x))$, the composition of the training map $f$ with a measurement $g$. The reference is $x_0 = \mathbf{1}$ (full unweighted training). The Taylor expansion around $x_0$ provides the interaction decomposition:
\[
  Q(x) = \sum_{r=0}^\infty T_r(x - x_0),
\]
where $T_r$ is the degree-$r$ homogeneous component. This decomposition is exact wherever $Q$ is smooth. \textbf{Diagnostic Question 1: Affirmative.}

\section{Stability: Diagnostic Question 2}

\begin{definition}[Gunn Stability \cite{Gunn2026}]
There exist $C, c > 0$ such that $\|T_r\| \leq Ce^{-cr}$ for all $r \geq 0$.
\end{definition}

This says higher-order interactions among training examples decay exponentially. Setting $\sigma_r = Ce^{-cr}$ gives a summable stability sequence. \textbf{Diagnostic Question 2: Affirmative under Gunn's stability assumption} (empirically supported for small models; open for frontier scale).

\section{Compression: Diagnostic Question 3}

Gunn's polynomial sketch \cite{Gunn2026} evaluates $Q(x_0 + \zeta\psi)$ for random complex directions $\psi$, extracting random projections of the derivative tensors via forward-mode automatic differentiation over truncated polynomial rings. The sketch achieves $(d, O(\varepsilon))$-sufficiency with sketch dimension $O(\log(1/\delta)/\varepsilon^2)$.  \textbf{Diagnostic Question 3: Affirmative.}

\section{Reconstruction and Alignment}

By Theorem~\ref{thm:reconstruction}:
\[
  E_{\mathrm{total}} \leq \frac{Ce^{-cd}}{1-e^{-c}} + O(\varepsilon) + \varepsilon_0.
\]
The sketch achieves \textbf{alignment by coincidence}: designed for computational efficiency, it happens to preserve derivative structure calibrated to the stability sequence.

\section{Diagnostic Audit}

\begin{center}\small
\begin{tabular}{lll}
\toprule
Question & Status & Evidence \\
\midrule
Decomposition exact? & Yes & Taylor series \\
Stability summable? & Yes (Gunn) & $\|T_r\| \leq Ce^{-cr}$, small-model tests \\
Compression sufficient? & Yes & Symmetric-subspace identity \\
Depths aligned? & Yes & $d_\pi = d_Q = d$ \\
Degeneracy avoided? & Yes & Well-separated sketch classes \\
\bottomrule
\end{tabular}
\end{center}

\section{Lesson}

Reconstruction works not because the training history is stored, but because the query depends on it in a compressible way. The stability assumption is the load-bearing hypothesis. Open obligation: verify exponential derivative decay at frontier-model scale.
